% =============================================================================
%  Part II: Cubic Scatterer  (Sections 11--16, 18)
% =============================================================================

\part*{Part II: Cubic Scatterer}

\section{Motivation: Why Cubes?}

The T-matrix for a cube is fundamentally more useful for effective medium
theory than the T-matrix for a sphere, because cubes tile three-dimensional
space.

A sphere is the natural scatterer for studying isolated inclusions in a
matrix, but when the goal is to model a heterogeneous medium as a
homogenised effective material, one must partition the entire volume into
non-overlapping cells.  Spheres cannot fill space without gaps: even the
densest packing of equal spheres leaves roughly 26\% void \citep{Hales2005}.
Cubes, by contrast, are space-filling: they tesselate $\mathbb{R}^3$ with
zero overlap and zero void.

This simple geometric fact has far-reaching consequences for effective
medium theory (see Section~17).

\subsection{Semi-Analytical Approach}

For the sphere, the radial decomposition of $G_{ij}$ and the angular
integral identities allow exact closed-form evaluation of all volume
integrals.  For the cube, no such radial-angular factorisation exists.
Instead, we use a semi-analytical approach:

\begin{enumerate}
\item Taylor-expand the Green's tensor $G_{ij}(\bx)$ about the centre of
  the cube ($\bx=\mathbf{0}$) as a polynomial in $(x_1,x_2,x_3)$, keeping
  terms through order $(\omega a)^7$.

\item Integrate each monomial $x_1^{p_1} x_2^{p_2} x_3^{p_3}$ over the
  cube $[-a,a]^3$ analytically using the factorisation:
\begin{equation}
  \int_{-a}^{a}\!\int_{-a}^{a}\!\int_{-a}^{a}
  x_1^{p_1} x_2^{p_2} x_3^{p_3}\,d^3x
  = \prod_{m=1}^{3} \frac{2\,a^{p_m+1}}{p_m+1}
  \qquad(\text{all } p_m \text{ even; zero otherwise}).
\label{eq:cubemono}
\end{equation}
\end{enumerate}
This gives the same leading-order behaviour as the sphere (the $O(ka)^3$
Rayleigh terms), plus a cubic-anisotropy correction at $O(ka)^7$ that has
no analogue for the sphere.


\section{Cubic Anisotropy Tensor}

The critical difference between sphere and cube appears in the fourth
moment of the scatterer volume.  For any centrosymmetric shape:
\begin{equation}
  \int_D x_i\,x_j\,x_k\,x_l\,d^3x
  = P\,(\delta_{ij}\delta_{kl} + \delta_{ik}\delta_{jl}
    + \delta_{il}\delta_{jk})
  + Q\,E_{ijkl},
\label{eq:4thmoment}
\end{equation}
where $E_{ijkl} = \sum_m \delta_{im}\delta_{jm}\delta_{km}\delta_{lm}$ is
the \emph{cubic anisotropy tensor} (equals~1 when all four indices are
equal, zero otherwise).

\textbf{Sphere:}  By spherical symmetry, $Q=0$.  The fourth moment is
purely isotropic.

\textbf{Cube} $[-a,a]^3$:  Direct computation gives
\begin{equation}
  \int x_1^4\,d^3x = \frac{8a^7}{5}, \qquad
  \int x_1^2 x_2^2\,d^3x = \frac{8a^7}{9},
\label{eq:cube4th}
\end{equation}
so $P = 8a^7/9$ and $Q = -16a^7/15 \neq 0$.  The cube has genuine cubic
anisotropy.


\section{Volume Integrals (Cube)}

\subsection{Green's Tensor Integral $\Gamma_0^{\text{cube}}$}

The zeroth-order integral
$\int_{\text{cube}} G_{ij}\,d^3x = \Gamma_0^{\text{cube}}\,\delta_{ij}$
is still isotropic (only the second moment enters, and
$\int x_i x_j\,d^3x \propto \delta_{ij}$ for the cube).  The result
differs from the sphere only through the geometry of the integration domain.

\subsection{Second-Derivative Integral: Three Parameters}

For the cube, the second-derivative integral acquires a third independent
tensor structure:
\begin{equation}
  \int_{\text{cube}} G_{ij,kl}\,d^3x
  = A^c\,\delta_{ij}\delta_{kl}
  + B^c\,(\delta_{ik}\delta_{jl} + \delta_{il}\delta_{jk})
  + C^c\,E_{ijkl},
\label{eq:ABCcube}
\end{equation}
where $C^c$ is the cubic anisotropy coefficient.  Computing the three
independent components:
\begin{equation}
  A^c = I_{1122}, \qquad
  B^c = I_{1212}, \qquad
  C^c = I_{1111} - I_{1122} - 2\,I_{1212}.
\label{eq:ABCdef}
\end{equation}
SymPy computes $C^c = O\bigl((\omega a)^7\bigr)$, confirming it arises
from the fourth-moment anisotropy of the cube.


\section{T-Matrix Coupling Coefficients (Cube)}

The T-matrix tensor for the cube has cubic (not isotropic) symmetry:
\begin{equation}
  T_{mnlp}^c = T_1^c\,\delta_{mn}\delta_{lp}
  + T_2^c\,(\delta_{ml}\delta_{np} + \delta_{mp}\delta_{nl})
  + T_3^c\,E_{mnlp}.
\label{eq:Tcube}
\end{equation}
The third coefficient $T_3^c$ arises from $C^c$ and vanishes for a sphere
($C=0 \Rightarrow T_3=0$).


\section{Self-Consistent Solution: The $6\times 6$ Voigt Matrix
  and Its Eigenvalue Decomposition}

\subsection{The Voigt Representation}

The self-consistent strain equation $\varepsilon_{mn} = \varepsilon_{mn}^0
+ T_{mnlp}\,\varepsilon_{lp}$ is a $3\times 3$ symmetric tensor equation.
In Voigt notation, the six independent strain components
\begin{equation}
  \boldsymbol{\varepsilon}_V
  = [\varepsilon_{11},\;\varepsilon_{22},\;\varepsilon_{33},\;
     2\varepsilon_{23},\;2\varepsilon_{13},\;2\varepsilon_{12}]^T
\label{eq:Voigt}
\end{equation}
are mapped to a 6-vector, and $T_{mnlp}$ becomes a $6\times 6$ matrix
$\bT_V$.

For the cubic T-tensor~\eqref{eq:Tcube}, $\bT_V$ has the block-diagonal
structure:
\begin{equation}
  \bT_V = \begin{pmatrix} \bD & \mathbf{0} \\ \mathbf{0} & \bS \end{pmatrix},
\label{eq:TVblock}
\end{equation}
where the $3\times 3$ diagonal-strain block is
\begin{equation}
  \bD = T_1\,\bJ + (2T_2 + T_3)\,\bI_3
  = \begin{pmatrix}
    T_1+2T_2+T_3 & T_1 & T_1 \\
    T_1 & T_1+2T_2+T_3 & T_1 \\
    T_1 & T_1 & T_1+2T_2+T_3
  \end{pmatrix}
\label{eq:Dblock}
\end{equation}
(with $\bJ$ the $3\times 3$ matrix of ones) and the $3\times 3$
off-diagonal shear block is
\begin{equation}
  \bS = 2T_2\,\bI_3.
\label{eq:Sblock}
\end{equation}

The physical reason for the block-diagonal structure is that the cubic
anisotropy tensor $E_{mnlp}$ only contributes when $m=n$ and $l=p$ (all
indices equal).  When either pair is off-diagonal ($m\neq n$ or $l\neq p$),
$E_{mnlp}=0$, so the $T_3$ term drops out and the off-diagonal shear
components decouple with the isotropic eigenvalue $2T_2$.

\subsection{Eigenvalue Decomposition: 4 Scalar Amplification Factors}

The self-consistent equation $\boldsymbol{\varepsilon} =
\boldsymbol{\varepsilon}^0 + \bT_V\,\boldsymbol{\varepsilon}$ has the
formal solution
$\boldsymbol{\varepsilon} = (\bI_6 - \bT_V)^{-1}\,\boldsymbol{\varepsilon}^0$.
Instead of inverting the full $6\times 6$ matrix symbolically (which
involves extremely expensive \texttt{simplify()} calls on nested rational
expressions), we exploit the block structure to reduce everything to four
scalar equations.

\textbf{Block $\bS$ (off-diagonal shear):}  Already diagonal.  The
amplification factor is
\begin{equation}
  A_e^{(\text{off})} = \frac{1}{1-2T_2}.
\label{eq:Aeoff}
\end{equation}

\textbf{Block $\bD$ (diagonal strains):}  The matrix
$\bD = T_1\,\bJ + (2T_2+T_3)\,\bI_3$ has the well-known eigenstructure:

\begin{itemize}
\item \textbf{Dilatation mode:}  eigenvector $[1,1,1]^T$ (uniform
  expansion), eigenvalue $3T_1 + 2T_2 + T_3$:
\begin{equation}
  A_\theta = \frac{1}{1-3T_1-2T_2-T_3}.
\label{eq:Athcube}
\end{equation}

\item \textbf{Diagonal deviatoric modes:}  eigenvectors $[1,-1,0]^T$ and
  $[1,1,-2]^T$ (traceless diagonal strain), eigenvalue $2T_2+T_3$:
\begin{equation}
  A_e^{(\text{diag})} = \frac{1}{1-2T_2-T_3}.
\label{eq:Aediag}
\end{equation}
\end{itemize}
Plus the displacement amplification (unchanged from the sphere):
\begin{equation}
  A_u = \frac{1}{1-\omega^2\Delta\rho\,\Gamma_0^{\text{cube}}}.
\label{eq:Aucube}
\end{equation}

For the sphere ($T_3=0$): $A_e^{(\text{off})} = A_e^{(\text{diag})} \equiv
A_e$ and we recover the 3 amplification factors of Section~8.


\section{Effective Contrasts (Cube): Scalar Derivation}

The effective stiffness contrast is
$\Delta c^*_{mnlp} = \Delta c_{mnuv}\,A_{uv,lp}$, where $A$ is the
amplification operator.  We avoid the full $6\times 6$ matrix product by
decomposing the incident strain into the three eigenspaces of~$A$ and
applying each scalar amplification factor separately.

The isotropic contrast $\Delta c_{mnuv}\,\varepsilon_{uv}^*
= \Delta\lambda\,\theta^*\,\delta_{mn} + 2\Delta\mu\,\varepsilon_{mn}^*$.
Substituting the amplified strain:
\begin{equation}
  \varepsilon_{mn}^*
  = \underbrace{\frac{A_\theta\,\theta^0}{3}\,\delta_{mn}}_{\text{dilatation}}
  + \underbrace{A_e^{(\text{diag})}\,e_{mn}^{0,\text{diag}}}_{\text{diag.\ deviatoric}}
  + \underbrace{A_e^{(\text{off})}\,e_{mn}^{0,\text{off}}}_{\text{off-diag.\ shear}}
\label{eq:ampstrain}
\end{equation}
gives, after algebra:

\medskip
\noindent\textbf{Effective density:}
\begin{equation}
  \Delta\rho^{*,c} = \frac{\Delta\rho}
  {1-\omega^2\Delta\rho\,\Gamma_0^{\text{cube}}}.
\label{eq:Drhostar_c}
\end{equation}

\noindent\textbf{Off-diagonal effective shear modulus:}
\begin{equation}
  \Delta\mu^{*,\text{off}} = \frac{\Delta\mu}{1-2T_2^c}.
\label{eq:Dmuoff}
\end{equation}

\noindent\textbf{Diagonal effective shear modulus:}
\begin{equation}
  \Delta\mu^{*,\text{diag}} = \frac{\Delta\mu}{1-2T_2^c-T_3^c}.
\label{eq:Dmudiag}
\end{equation}

\noindent\textbf{Effective Lam\'e parameter:}
\begin{equation}
  \Delta\lambda^{*,c}
  = \frac{\Delta\lambda + \tfrac{2}{3}\Delta\mu}
         {1-3T_1^c-2T_2^c-T_3^c}
  - \frac{\Delta\mu}{3}
    \left(\frac{1}{1-2T_2^c-T_3^c} + \frac{1}{1-2T_2^c}\right).
\label{eq:Dlamstar_c}
\end{equation}

The cubic anisotropy of the effective shear modulus is:
\begin{equation}
  \Delta\mu^{*,\text{diag}} - \Delta\mu^{*,\text{off}}
  = \frac{\Delta\mu\,T_3^c}
  {(1-2T_2^c)(1-2T_2^c-T_3^c)}
  = O\bigl((\omega a)^7\bigr).
\label{eq:cubicaniso}
\end{equation}
